%La lectura por doce agrupaciones es la columna vertebral del análisis sectorial. Dentro de cada subsección se incluye, cuando corresponde, un zoom a 25 agrupaciones para mostrar qué ocurre al abrir la actividad amplia en subsectores más específicos. En las actividades donde la apertura a 25 agrupaciones coincide con la de 12, pero la apertura a 61 sí ofrece mayor detalle, el zoom usa esa desagregación más fina.
\textbf{A continuación se presenta la descomposición del crecimiento de cada una de las doce agrupaciones de actividad económica CIIU.} En particular, se presenta el PIB real de la actividad, el número de ocupados, el PIB por trabajador, las horas semanales promedio por trabajador y el PIB por hora trabajada para los años 2010 y 2025. La lectura conjunta de estas variables permite distinguir si los cambios de productividad responden principalmente al dinamismo del producto, a variaciones en el empleo, a cambios en las horas trabajadas o a una combinación de estos factores. Las descripciones siguen la agregación usada por el DANE; por eso, en algunos casos reúnen actividades económicas muy distintas dentro de una misma agrupación.

\subsection{Agropecuario}

Esta agrupación reúne agricultura, ganadería, caza, silvicultura y pesca.

\begin{table}[H]
\centering
\caption{Agropecuario: PIB, ocupados, horas y productividad laboral, 2010--2025}
\label{tab:sector_a_productividad}
\small
\begin{tabular}{lrrr}
\toprule
Indicador & 2010 & 2025 & Crec. anual \\
\midrule
PIB real  (Billones de pesos de 2015) & 39,9 & 63,9 & 3,18\% \\
Ocupados (Millones) & 3,0 & 3,1 & 0,38\% \\
PIB por trabajador (Millones de pesos de 2015) & 13,5 & 20,4 & 2,79\% \\
Horas semanales por trabajador & 45,0 & 42,0 & -0,46\% \\
PIB por hora trabajada (Miles de pesos de 2015) & 5,8 & 9,4 & 3,27\% \\
\bottomrule
\end{tabular}
\end{table}

\textbf{Frente al agregado nacional, la comparación variable por variable muestra diferencias relevantes.} El PIB real de la actividad registró una tasa anualizada de 3,18\%, muy cerca del agregado (3,16\%). La ocupación en la actividad registró una tasa anualizada de 0,38\%, por debajo del agregado (2,09\%). El PIB por trabajador registró una tasa anualizada de 2,79\%, por encima del agregado (1,05\%). Las horas semanales por trabajador registraron una tasa anualizada de -0,46\%, muy cerca de la caída agregada (-0,45\%). El PIB por hora trabajada registró una tasa anualizada de 3,27\%, por encima del agregado (1,51\%).

\textbf{En niveles de 2025, el tamaño relativo y la productividad de la actividad económica también importan.} La actividad representó 6,3\% del PIB real agregado, con 63,9 billones de pesos de 2015, y concentró 13,6\% de los ocupados, con 3,1 millones de personas. El PIB por trabajador fue 20,4 millones de pesos de 2015 por ocupado, por debajo del agregado (44,5 millones). Las horas semanales por trabajador fueron 42,0, muy cerca del agregado (43,7). El PIB por hora trabajada fue 9,4 mil pesos de 2015 por hora, por debajo del agregado (19,6 mil).

\textbf{Zoom a 61 agrupaciones.} En agropecuario, la apertura a 25 agrupaciones coincide con la agrupación de 12. Por eso, el zoom usa la apertura a 61 agrupaciones, que permite mirar subactividades más específicas.

\begin{table}[H]
\centering
\caption{Agropecuario: apertura de productividad laboral a 61 agrupaciones, 2010--2025}
\label{tab:sector_a_zoom61}
\scriptsize
\begin{tabular}{p{0.40\textwidth}rrrr}
\toprule
Actividad económica & PIB/trab. 2025 & Crec. & PIB/hora 2025 & Crec. \\
\midrule
Silvicultura & 70,2 & 5,15\% & 32,9 & 6,27\% \\
Pesca y acuicultura & 24,6 & 4,79\% & 10,9 & 5,75\% \\
Cultivos y apoyo agropecuario; Café; Ganadería & 19,9 & 2,74\% & 9,1 & 3,19\% \\
\bottomrule
\end{tabular}
\caption*{\footnotesize Nota: PIB por trabajador en millones de pesos constantes de 2015; PIB por hora en miles de pesos constantes de 2015. Cuando la GEIH no separa ocupados y horas al mismo nivel de las cuentas nacionales, se agrupan las subactividades del DANE hasta el nivel laboral comparable. Fuente: cálculos propios con DANE y GEIH.}
\end{table}

En esta apertura, el mayor crecimiento del PIB por trabajador se observa en Silvicultura (5,15\%), mientras que el menor se registra en Cultivos y apoyo agropecuario; Café; Ganadería (2,74\%). Por hora trabajada, el mejor desempeño corresponde a Silvicultura (6,27\%) y el más rezagado a Cultivos y apoyo agropecuario; Café; Ganadería (3,19\%).

\subsection{Minas}

Esta agrupación reúne explotación de minas y canteras. En la CIIU Rev. 4 A.C. corresponde a B; divisiones 05--09.

\begin{table}[H]
\centering
\caption{Minas: PIB, ocupados, horas y productividad laboral, 2010--2025}
\label{tab:sector_b_productividad}
\small
\begin{tabular}{lrrr}
\toprule
Indicador & 2010 & 2025 & Crec. anual \\
\midrule
PIB real  (Billones de pesos de 2015) & 38,4 & 34,7 & -0,67\% \\
Ocupados (Millones) & 0,2 & 0,3 & 3,04\% \\
PIB por trabajador (Millones de pesos de 2015) & 200,8 & 115,8 & -3,60\% \\
Horas semanales por trabajador & 49,7 & 46,0 & -0,52\% \\
PIB por hora trabajada (Miles de pesos de 2015) & 77,6 & 48,4 & -3,09\% \\
\bottomrule
\end{tabular}
\end{table}

\textbf{Frente al agregado nacional, la comparación variable por variable muestra diferencias relevantes.} El PIB real de la actividad registró una tasa anualizada de -0,67\%, por debajo del agregado (3,16\%). La ocupación en la actividad registró una tasa anualizada de 3,04\%, por encima del agregado (2,09\%). El PIB por trabajador registró una tasa anualizada de -3,60\%, por debajo del agregado (1,05\%). Las horas semanales por trabajador registraron una tasa anualizada de -0,52\%, muy cerca de la caída agregada (-0,45\%). El PIB por hora trabajada registró una tasa anualizada de -3,09\%, por debajo del agregado (1,51\%).

\textbf{En niveles de 2025, el tamaño relativo y la productividad de la actividad económica también importan.} La actividad representó 3,4\% del PIB real agregado, con 34,7 billones de pesos de 2015, y concentró 1,3\% de los ocupados, con 0,3 millones de personas. El PIB por trabajador fue 115,8 millones de pesos de 2015 por ocupado, por encima del agregado (44,5 millones). Las horas semanales por trabajador fueron 46,0, por encima del agregado (43,7). El PIB por hora trabajada fue 48,4 mil pesos de 2015 por hora, por encima del agregado (19,6 mil).

\textbf{Zoom a 61 agrupaciones.} En minas, la apertura a 25 agrupaciones coincide con la agrupación de 12. Por eso, el zoom usa la apertura a 61 agrupaciones, que permite mirar subactividades más específicas.

\begin{table}[H]
\centering
\caption{Minas: apertura de productividad laboral a 61 agrupaciones, 2010--2025}
\label{tab:sector_b_zoom61}
\scriptsize
\begin{tabular}{p{0.40\textwidth}rrrr}
\toprule
Actividad económica & PIB/trab. 2025 & Crec. & PIB/hora 2025 & Crec. \\
\midrule
Carbón; Petróleo, gas y apoyo conexo; Apoyo a otras actividades mineras & 308,0 & -2,44\% & 124,1 & -1,90\% \\
Minerales metalíferos; Otras minas y canteras & 22,1 & -5,13\% & 9,4 & -4,68\% \\
\bottomrule
\end{tabular}
\caption*{\footnotesize Nota: PIB por trabajador en millones de pesos constantes de 2015; PIB por hora en miles de pesos constantes de 2015. Cuando la GEIH no separa ocupados y horas al mismo nivel de las cuentas nacionales, se agrupan las subactividades del DANE hasta el nivel laboral comparable. Fuente: cálculos propios con DANE y GEIH.}
\end{table}

En esta apertura, el mayor crecimiento del PIB por trabajador se observa en Carbón; Petróleo, gas y apoyo conexo; Apoyo a otras actividades mineras (-2,44\%), mientras que el menor se registra en Minerales metalíferos; Otras minas y canteras (-5,13\%). Por hora trabajada, el mejor desempeño corresponde a Carbón; Petróleo, gas y apoyo conexo; Apoyo a otras actividades mineras (-1,90\%) y el más rezagado a Minerales metalíferos; Otras minas y canteras (-4,68\%).

\subsection{Manufactura}

Esta agrupación reúne industrias manufactureras. En la CIIU Rev. 4 A.C. corresponde a C; divisiones 10--33.

\begin{table}[H]
\centering
\caption{Manufactura: PIB, ocupados, horas y productividad laboral, 2010--2025}
\label{tab:sector_c_productividad}
\small
\begin{tabular}{lrrr}
\toprule
Indicador & 2010 & 2025 & Crec. anual \\
\midrule
PIB real  (Billones de pesos de 2015) & 88,0 & 113,1 & 1,69\% \\
Ocupados (Millones) & 2,1 & 2,5 & 1,07\% \\
PIB por trabajador (Millones de pesos de 2015) & 41,1 & 45,0 & 0,61\% \\
Horas semanales por trabajador & 46,6 & 44,3 & -0,34\% \\
PIB por hora trabajada (Miles de pesos de 2015) & 16,9 & 19,5 & 0,95\% \\
\bottomrule
\end{tabular}
\end{table}

\textbf{Frente al agregado nacional, la comparación variable por variable muestra diferencias relevantes.} El PIB real de la actividad registró una tasa anualizada de 1,69\%, por debajo del agregado (3,16\%). La ocupación en la actividad registró una tasa anualizada de 1,07\%, por debajo del agregado (2,09\%). El PIB por trabajador registró una tasa anualizada de 0,61\%, por debajo del agregado (1,05\%). Las horas semanales por trabajador registraron una tasa anualizada de -0,34\%, muy cerca de la caída agregada (-0,45\%). El PIB por hora trabajada registró una tasa anualizada de 0,95\%, por debajo del agregado (1,51\%).

\textbf{En niveles de 2025, el tamaño relativo y la productividad de la actividad económica también importan.} La actividad representó 11,1\% del PIB real agregado, con 113,1 billones de pesos de 2015, y concentró 10,9\% de los ocupados, con 2,5 millones de personas. El PIB por trabajador fue 45,0 millones de pesos de 2015 por ocupado, muy cerca del agregado (44,5 millones). Las horas semanales por trabajador fueron 44,3, muy cerca del agregado (43,7). El PIB por hora trabajada fue 19,5 mil pesos de 2015 por hora, muy cerca del agregado (19,6 mil).

\textbf{Zoom a 25 agrupaciones.} Dentro de manufactura, la apertura a 25 agrupaciones permite ver diferencias que el promedio de la actividad oculta.

\begin{table}[H]
\centering
\caption{Manufactura: apertura de productividad laboral a 25 agrupaciones, 2010--2025}
\label{tab:sector_c_zoom25}
\scriptsize
\begin{tabular}{p{0.40\textwidth}rrrr}
\toprule
Actividad económica & PIB/trab. 2025 & Crec. & PIB/hora 2025 & Crec. \\
\midrule
Alimentos, bebidas y tabaco & 50,8 & 1,60\% & 21,8 & 1,70\% \\
Madera, papel e impresión & 46,9 & 1,46\% & 21,8 & 2,21\% \\
Textiles, confecciones y cuero & 17,5 & 0,74\% & 7,8 & 1,19\% \\
Metalurgia, maquinaria y equipo & 48,8 & 0,53\% & 20,5 & 1,00\% \\
Refinación, químicos y minerales & 106,6 & -0,75\% & 44,8 & -0,30\% \\
Muebles y otras manufactureras & 17,4 & -1,31\% & 7,6 & -1,06\% \\
\bottomrule
\end{tabular}
\caption*{\footnotesize Nota: PIB por trabajador en millones de pesos constantes de 2015; PIB por hora en miles de pesos constantes de 2015. Fuente: cálculos propios con DANE y GEIH.}
\end{table}

En esta apertura, el mayor crecimiento del PIB por trabajador se observa en Alimentos, bebidas y tabaco (1,60\%), mientras que el menor se registra en Muebles y otras manufactureras (-1,31\%). Por hora trabajada, el mejor desempeño corresponde a Madera, papel e impresión (2,21\%) y el más rezagado a Muebles y otras manufactureras (-1,06\%).

\subsection{Servicios públicos}

Esta agrupación reúne suministro de electricidad, gas, vapor y aire acondicionado; distribución de agua; evacuación y tratamiento de aguas residuales, gestión de desechos y actividades de saneamiento ambiental. En la CIIU Rev. 4 A.C. corresponde a D y E; divisiones 35--39.

\begin{table}[H]
\centering
\caption{Servicios públicos: PIB, ocupados, horas y productividad laboral, 2010--2025}
\label{tab:sector_d_e_productividad}
\small
\begin{tabular}{lrrr}
\toprule
Indicador & 2010 & 2025 & Crec. anual \\
\midrule
PIB real  (Billones de pesos de 2015) & 21,9 & 30,3 & 2,18\% \\
Ocupados (Millones) & 0,1 & 0,3 & 6,30\% \\
PIB por trabajador (Millones de pesos de 2015) & 180,1 & 99,5 & -3,88\% \\
Horas semanales por trabajador & 48,6 & 44,0 & -0,65\% \\
PIB por hora trabajada (Miles de pesos de 2015) & 71,3 & 43,5 & -3,25\% \\
\bottomrule
\end{tabular}
\end{table}

\textbf{Frente al agregado nacional, la comparación variable por variable muestra diferencias relevantes.} El PIB real de la actividad registró una tasa anualizada de 2,18\%, por debajo del agregado (3,16\%). La ocupación en la actividad registró una tasa anualizada de 6,30\%, por encima del agregado (2,09\%). El PIB por trabajador registró una tasa anualizada de -3,88\%, por debajo del agregado (1,05\%). Las horas semanales por trabajador registraron una tasa anualizada de -0,65\%, una caída más pronunciada que la del agregado (-0,45\%). El PIB por hora trabajada registró una tasa anualizada de -3,25\%, por debajo del agregado (1,51\%).

\textbf{En niveles de 2025, el tamaño relativo y la productividad de la actividad económica también importan.} La actividad representó 3,0\% del PIB real agregado, con 30,3 billones de pesos de 2015, y concentró 1,3\% de los ocupados, con 0,3 millones de personas. El PIB por trabajador fue 99,5 millones de pesos de 2015 por ocupado, por encima del agregado (44,5 millones). Las horas semanales por trabajador fueron 44,0, muy cerca del agregado (43,7). El PIB por hora trabajada fue 43,5 mil pesos de 2015 por hora, por encima del agregado (19,6 mil).

\textbf{Zoom a 25 agrupaciones.} Dentro de servicios públicos, la apertura a 25 agrupaciones permite ver diferencias que el promedio de la actividad oculta.

\begin{table}[H]
\centering
\caption{Servicios públicos: apertura de productividad laboral a 25 agrupaciones, 2010--2025}
\label{tab:sector_d_e_zoom25}
\scriptsize
\begin{tabular}{p{0.40\textwidth}rrrr}
\toprule
Actividad económica & PIB/trab. 2025 & Crec. & PIB/hora 2025 & Crec. \\
\midrule
Electricidad, gas y vapor & 262,2 & -0,13\% & 109,6 & 0,28\% \\
Agua, saneamiento y desechos & 37,9 & -6,64\% & 16,9 & -5,96\% \\
\bottomrule
\end{tabular}
\caption*{\footnotesize Nota: PIB por trabajador en millones de pesos constantes de 2015; PIB por hora en miles de pesos constantes de 2015. Fuente: cálculos propios con DANE y GEIH.}
\end{table}

En esta apertura, el mayor crecimiento del PIB por trabajador se observa en Electricidad, gas y vapor (-0,13\%), mientras que el menor se registra en Agua, saneamiento y desechos (-6,64\%). Por hora trabajada, el mejor desempeño corresponde a Electricidad, gas y vapor (0,28\%) y el más rezagado a Agua, saneamiento y desechos (-5,96\%).

\subsection{Construcción}

Esta agrupación reúne construcción. En la CIIU Rev. 4 A.C. corresponde a F; divisiones 41--43.

\begin{table}[H]
\centering
\caption{Construcción: PIB, ocupados, horas y productividad laboral, 2010--2025}
\label{tab:sector_f_productividad}
\small
\begin{tabular}{lrrr}
\toprule
Indicador & 2010 & 2025 & Crec. anual \\
\midrule
PIB real  (Billones de pesos de 2015) & 40,0 & 41,3 & 0,21\% \\
Ocupados (Millones) & 1,0 & 1,6 & 3,39\% \\
PIB por trabajador (Millones de pesos de 2015) & 41,8 & 26,1 & -3,08\% \\
Horas semanales por trabajador & 48,9 & 45,7 & -0,46\% \\
PIB por hora trabajada (Miles de pesos de 2015) & 16,4 & 11,0 & -2,63\% \\
\bottomrule
\end{tabular}
\end{table}

\textbf{Frente al agregado nacional, la comparación variable por variable muestra diferencias relevantes.} El PIB real de la actividad registró una tasa anualizada de 0,21\%, por debajo del agregado (3,16\%). La ocupación en la actividad registró una tasa anualizada de 3,39\%, por encima del agregado (2,09\%). El PIB por trabajador registró una tasa anualizada de -3,08\%, por debajo del agregado (1,05\%). Las horas semanales por trabajador registraron una tasa anualizada de -0,46\%, muy cerca de la caída agregada (-0,45\%). El PIB por hora trabajada registró una tasa anualizada de -2,63\%, por debajo del agregado (1,51\%).

\textbf{En niveles de 2025, el tamaño relativo y la productividad de la actividad económica también importan.} La actividad representó 4,0\% del PIB real agregado, con 41,3 billones de pesos de 2015, y concentró 6,9\% de los ocupados, con 1,6 millones de personas. El PIB por trabajador fue 26,1 millones de pesos de 2015 por ocupado, por debajo del agregado (44,5 millones). Las horas semanales por trabajador fueron 45,7, muy cerca del agregado (43,7). El PIB por hora trabajada fue 11,0 mil pesos de 2015 por hora, por debajo del agregado (19,6 mil).

\textbf{Zoom a 25 agrupaciones.} El DANE abre construcción en edificaciones, obras civiles y actividades especializadas de construcción. Sin embargo, la GEIH histórica comparable no permite separar ocupados y horas entre esas tres subactividades. Por esa razón, no se presenta un cuadro de PIB por trabajador o PIB por hora a este nivel. La medición comparable corresponde al total de construcción.

\subsection{Comercio, transporte y alojamiento}

Esta agrupación reúne comercio al por mayor y al por menor; reparación de vehículos automotores y motocicletas; transporte y almacenamiento; alojamiento y servicios de comida. En la CIIU Rev. 4 A.C. corresponde a G, H e I; divisiones 45--56.

\begin{table}[H]
\centering
\caption{Comercio, transporte y alojamiento: PIB, ocupados, horas y productividad laboral, 2010--2025}
\label{tab:sector_g_h_i_productividad}
\small
\begin{tabular}{lrrr}
\toprule
Indicador & 2010 & 2025 & Crec. anual \\
\midrule
PIB real  (Billones de pesos de 2015) & 107,6 & 179,1 & 3,45\% \\
Ocupados (Millones) & 5,4 & 7,4 & 2,14\% \\
PIB por trabajador (Millones de pesos de 2015) & 19,9 & 24,1 & 1,28\% \\
Horas semanales por trabajador & 50,2 & 45,8 & -0,61\% \\
PIB por hora trabajada (Miles de pesos de 2015) & 7,6 & 10,1 & 1,90\% \\
\bottomrule
\end{tabular}
\end{table}

\textbf{Frente al agregado nacional, la comparación variable por variable muestra diferencias relevantes.} El PIB real de la actividad registró una tasa anualizada de 3,45\%, por encima del agregado (3,16\%). La ocupación en la actividad registró una tasa anualizada de 2,14\%, muy cerca del agregado (2,09\%). El PIB por trabajador registró una tasa anualizada de 1,28\%, por encima del agregado (1,05\%). Las horas semanales por trabajador registraron una tasa anualizada de -0,61\%, muy cerca de la caída agregada (-0,45\%). El PIB por hora trabajada registró una tasa anualizada de 1,90\%, por encima del agregado (1,51\%).

\textbf{En niveles de 2025, el tamaño relativo y la productividad de la actividad económica también importan.} La actividad representó 17,5\% del PIB real agregado, con 179,1 billones de pesos de 2015, y concentró 32,4\% de los ocupados, con 7,4 millones de personas. El PIB por trabajador fue 24,1 millones de pesos de 2015 por ocupado, por debajo del agregado (44,5 millones). Las horas semanales por trabajador fueron 45,8, muy cerca del agregado (43,7). El PIB por hora trabajada fue 10,1 mil pesos de 2015 por hora, por debajo del agregado (19,6 mil).

\textbf{Zoom a 25 agrupaciones.} Dentro de comercio, transporte y alojamiento, la apertura a 25 agrupaciones permite ver diferencias que el promedio de la actividad oculta.

\begin{table}[H]
\centering
\caption{Comercio, transporte y alojamiento: apertura de productividad laboral a 25 agrupaciones, 2010--2025}
\label{tab:sector_g_h_i_zoom25}
\scriptsize
\begin{tabular}{p{0.40\textwidth}rrrr}
\toprule
Actividad económica & PIB/trab. 2025 & Crec. & PIB/hora 2025 & Crec. \\
\midrule
Comercio y reparación & 22,8 & 2,75\% & 9,5 & 3,16\% \\
Transporte y almacenamiento & 27,8 & -0,04\% & 10,9 & 0,94\% \\
Alojamiento y comida & 23,4 & -1,08\% & 10,7 & -0,42\% \\
\bottomrule
\end{tabular}
\caption*{\footnotesize Nota: PIB por trabajador en millones de pesos constantes de 2015; PIB por hora en miles de pesos constantes de 2015. Fuente: cálculos propios con DANE y GEIH.}
\end{table}

En esta apertura, el mayor crecimiento del PIB por trabajador se observa en Comercio y reparación (2,75\%), mientras que el menor se registra en Alojamiento y comida (-1,08\%). Por hora trabajada, el mejor desempeño corresponde a Comercio y reparación (3,16\%) y el más rezagado a Alojamiento y comida (-0,42\%).

\subsection{Información y comunicaciones}

Esta agrupación reúne información y comunicaciones. En la CIIU Rev. 4 A.C. corresponde a J; divisiones 58--63.

\begin{table}[H]
\centering
\caption{Información y comunicaciones: PIB, ocupados, horas y productividad laboral, 2010--2025}
\label{tab:sector_j_productividad}
\small
\begin{tabular}{lrrr}
\toprule
Indicador & 2010 & 2025 & Crec. anual \\
\midrule
PIB real  (Billones de pesos de 2015) & 18,3 & 31,3 & 3,66\% \\
Ocupados (Millones) & 0,4 & 0,4 & -0,08\% \\
PIB por trabajador (Millones de pesos de 2015) & 45,3 & 78,6 & 3,74\% \\
Horas semanales por trabajador & 49,6 & 44,4 & -0,73\% \\
PIB por hora trabajada (Miles de pesos de 2015) & 17,6 & 34,0 & 4,50\% \\
\bottomrule
\end{tabular}
\end{table}

\textbf{Frente al agregado nacional, la comparación variable por variable muestra diferencias relevantes.} El PIB real de la actividad registró una tasa anualizada de 3,66\%, por encima del agregado (3,16\%). La ocupación en la actividad registró una tasa anualizada de -0,08\%, por debajo del agregado (2,09\%). El PIB por trabajador registró una tasa anualizada de 3,74\%, por encima del agregado (1,05\%). Las horas semanales por trabajador registraron una tasa anualizada de -0,73\%, una caída más pronunciada que la del agregado (-0,45\%). El PIB por hora trabajada registró una tasa anualizada de 4,50\%, por encima del agregado (1,51\%).

\textbf{En niveles de 2025, el tamaño relativo y la productividad de la actividad económica también importan.} La actividad representó 3,1\% del PIB real agregado, con 31,3 billones de pesos de 2015, y concentró 1,7\% de los ocupados, con 0,4 millones de personas. El PIB por trabajador fue 78,6 millones de pesos de 2015 por ocupado, por encima del agregado (44,5 millones). Las horas semanales por trabajador fueron 44,4, muy cerca del agregado (43,7). El PIB por hora trabajada fue 34,0 mil pesos de 2015 por hora, por encima del agregado (19,6 mil).

\textbf{Zoom a 25 agrupaciones.} En esta actividad, la apertura a 25 agrupaciones coincide con la agrupación de 12: Información y comunicaciones. Por eso, el cuadro anterior ya resume la lectura relevante a este nivel.

\subsection{Financieras}

Esta agrupación reúne actividades financieras y de seguros. En la CIIU Rev. 4 A.C. corresponde a K; divisiones 64--66.

\begin{table}[H]
\centering
\caption{Financieras: PIB, ocupados, horas y productividad laboral, 2010--2025}
\label{tab:sector_k_productividad}
\small
\begin{tabular}{lrrr}
\toprule
Indicador & 2010 & 2025 & Crec. anual \\
\midrule
PIB real  (Billones de pesos de 2015) & 22,3 & 53,1 & 5,95\% \\
Ocupados (Millones) & 0,2 & 0,4 & 4,49\% \\
PIB por trabajador (Millones de pesos de 2015) & 101,0 & 124,5 & 1,40\% \\
Horas semanales por trabajador & 45,2 & 44,2 & -0,15\% \\
PIB por hora trabajada (Miles de pesos de 2015) & 43,0 & 54,2 & 1,56\% \\
\bottomrule
\end{tabular}
\end{table}

\textbf{Frente al agregado nacional, la comparación variable por variable muestra diferencias relevantes.} El PIB real de la actividad registró una tasa anualizada de 5,95\%, por encima del agregado (3,16\%). La ocupación en la actividad registró una tasa anualizada de 4,49\%, por encima del agregado (2,09\%). El PIB por trabajador registró una tasa anualizada de 1,40\%, por encima del agregado (1,05\%). Las horas semanales por trabajador registraron una tasa anualizada de -0,15\%, una caída menos pronunciada que la del agregado (-0,45\%). El PIB por hora trabajada registró una tasa anualizada de 1,56\%, muy cerca del agregado (1,51\%).

\textbf{En niveles de 2025, el tamaño relativo y la productividad de la actividad económica también importan.} La actividad representó 5,2\% del PIB real agregado, con 53,1 billones de pesos de 2015, y concentró 1,9\% de los ocupados, con 0,4 millones de personas. El PIB por trabajador fue 124,5 millones de pesos de 2015 por ocupado, por encima del agregado (44,5 millones). Las horas semanales por trabajador fueron 44,2, muy cerca del agregado (43,7). El PIB por hora trabajada fue 54,2 mil pesos de 2015 por hora, por encima del agregado (19,6 mil).

\textbf{Zoom a 25 agrupaciones.} En esta actividad, la apertura a 25 agrupaciones coincide con la agrupación de 12: Financieras y seguros. Por eso, el cuadro anterior ya resume la lectura relevante a este nivel.

\subsection{Inmobiliarias}

Esta agrupación reúne actividades inmobiliarias. En la CIIU Rev. 4 A.C. corresponde a L; división 68.

\begin{table}[H]
\centering
\caption{Inmobiliarias: PIB, ocupados, horas y productividad laboral, 2010--2025}
\label{tab:sector_l_productividad}
\small
\begin{tabular}{lrrr}
\toprule
Indicador & 2010 & 2025 & Crec. anual \\
\midrule
PIB real  (Billones de pesos de 2015) & 59,9 & 90,1 & 2,75\% \\
Ocupados (Millones) & 0,2 & 0,3 & 2,64\% \\
PIB por trabajador (Millones de pesos de 2015) & 280,0 & 284,6 & 0,11\% \\
Horas semanales por trabajador & 51,1 & 45,8 & -0,73\% \\
PIB por hora trabajada (Miles de pesos de 2015) & 105,3 & 119,5 & 0,85\% \\
\bottomrule
\end{tabular}
\end{table}

\textbf{Frente al agregado nacional, la comparación variable por variable muestra diferencias relevantes.} El PIB real de la actividad registró una tasa anualizada de 2,75\%, por debajo del agregado (3,16\%). La ocupación en la actividad registró una tasa anualizada de 2,64\%, por encima del agregado (2,09\%). El PIB por trabajador registró una tasa anualizada de 0,11\%, por debajo del agregado (1,05\%). Las horas semanales por trabajador registraron una tasa anualizada de -0,73\%, una caída más pronunciada que la del agregado (-0,45\%). El PIB por hora trabajada registró una tasa anualizada de 0,85\%, por debajo del agregado (1,51\%).

\textbf{En niveles de 2025, el tamaño relativo y la productividad de la actividad económica también importan.} La actividad representó 8,8\% del PIB real agregado, con 90,1 billones de pesos de 2015, y concentró 1,4\% de los ocupados, con 0,3 millones de personas. El PIB por trabajador fue 284,6 millones de pesos de 2015 por ocupado, por encima del agregado (44,5 millones). Las horas semanales por trabajador fueron 45,8, muy cerca del agregado (43,7). El PIB por hora trabajada fue 119,5 mil pesos de 2015 por hora, por encima del agregado (19,6 mil).

\textbf{Zoom a 25 agrupaciones.} En esta actividad, la apertura a 25 agrupaciones coincide con la agrupación de 12: Inmobiliarias. Por eso, el cuadro anterior ya resume la lectura relevante a este nivel.

\subsection{Profesionales y administrativas}

Esta agrupación reúne actividades profesionales, científicas y técnicas; actividades de servicios administrativos y de apoyo. En la CIIU Rev. 4 A.C. corresponde a M y N; divisiones 69--82.

\begin{table}[H]
\centering
\caption{Profesionales y administrativas: PIB, ocupados, horas y productividad laboral, 2010--2025}
\label{tab:sector_m_n_productividad}
\small
\begin{tabular}{lrrr}
\toprule
Indicador & 2010 & 2025 & Crec. anual \\
\midrule
PIB real  (Billones de pesos de 2015) & 45,4 & 70,1 & 2,94\% \\
Ocupados (Millones) & 0,8 & 1,7 & 4,94\% \\
PIB por trabajador (Millones de pesos de 2015) & 56,6 & 42,4 & -1,90\% \\
Horas semanales por trabajador & 41,9 & 39,6 & -0,38\% \\
PIB por hora trabajada (Miles de pesos de 2015) & 26,0 & 20,6 & -1,52\% \\
\bottomrule
\end{tabular}
\end{table}

\textbf{Frente al agregado nacional, la comparación variable por variable muestra diferencias relevantes.} El PIB real de la actividad registró una tasa anualizada de 2,94\%, por debajo del agregado (3,16\%). La ocupación en la actividad registró una tasa anualizada de 4,94\%, por encima del agregado (2,09\%). El PIB por trabajador registró una tasa anualizada de -1,90\%, por debajo del agregado (1,05\%). Las horas semanales por trabajador registraron una tasa anualizada de -0,38\%, muy cerca de la caída agregada (-0,45\%). El PIB por hora trabajada registró una tasa anualizada de -1,52\%, por debajo del agregado (1,51\%).

\textbf{En niveles de 2025, el tamaño relativo y la productividad de la actividad económica también importan.} La actividad representó 6,9\% del PIB real agregado, con 70,1 billones de pesos de 2015, y concentró 7,2\% de los ocupados, con 1,7 millones de personas. El PIB por trabajador fue 42,4 millones de pesos de 2015 por ocupado, muy cerca del agregado (44,5 millones). Las horas semanales por trabajador fueron 39,6, por debajo del agregado (43,7). El PIB por hora trabajada fue 20,6 mil pesos de 2015 por hora, por encima del agregado (19,6 mil).

\textbf{Zoom a 61 agrupaciones.} En profesionales y administrativas, la apertura a 25 agrupaciones coincide con la agrupación de 12. Por eso, el zoom usa la apertura a 61 agrupaciones, que permite mirar subactividades más específicas.

\begin{table}[H]
\centering
\caption{Profesionales y administrativas: apertura de productividad laboral a 61 agrupaciones, 2010--2025}
\label{tab:sector_m_n_zoom61}
\scriptsize
\begin{tabular}{p{0.40\textwidth}rrrr}
\toprule
Actividad económica & PIB/trab. 2025 & Crec. & PIB/hora 2025 & Crec. \\
\midrule
Servicios administrativos y de apoyo & 36,0 & -1,61\% & 18,0 & -1,20\% \\
Profesionales, científicas y técnicas & 54,7 & -2,16\% & 25,2 & -1,86\% \\
\bottomrule
\end{tabular}
\caption*{\footnotesize Nota: PIB por trabajador en millones de pesos constantes de 2015; PIB por hora en miles de pesos constantes de 2015. Cuando la GEIH no separa ocupados y horas al mismo nivel de las cuentas nacionales, se agrupan las subactividades del DANE hasta el nivel laboral comparable. Fuente: cálculos propios con DANE y GEIH.}
\end{table}

En esta apertura, el mayor crecimiento del PIB por trabajador se observa en Servicios administrativos y de apoyo (-1,61\%), mientras que el menor se registra en Profesionales, científicas y técnicas (-2,16\%). Por hora trabajada, el mejor desempeño corresponde a Servicios administrativos y de apoyo (-1,20\%) y el más rezagado a Profesionales, científicas y técnicas (-1,86\%).

\subsection{Adm. pública, educación y salud}

Esta agrupación reúne administración pública y defensa; planes de seguridad social de afiliación obligatoria; educación; actividades de atención de la salud humana y de servicios sociales. En la CIIU Rev. 4 A.C. corresponde a O, P y Q; divisiones 84--88.

\begin{table}[H]
\centering
\caption{Adm. pública, educación y salud: PIB, ocupados, horas y productividad laboral, 2010--2025}
\label{tab:sector_o_p_q_productividad}
\small
\begin{tabular}{lrrr}
\toprule
Indicador & 2010 & 2025 & Crec. anual \\
\midrule
PIB real  (Billones de pesos de 2015) & 85,4 & 165,3 & 4,51\% \\
Ocupados (Millones) & 2,0 & 2,8 & 2,50\% \\
PIB por trabajador (Millones de pesos de 2015) & 43,6 & 58,4 & 1,96\% \\
Horas semanales por trabajador & 43,3 & 43,0 & -0,04\% \\
PIB por hora trabajada (Miles de pesos de 2015) & 19,4 & 26,1 & 2,00\% \\
\bottomrule
\end{tabular}
\end{table}

\textbf{Frente al agregado nacional, la comparación variable por variable muestra diferencias relevantes.} El PIB real de la actividad registró una tasa anualizada de 4,51\%, por encima del agregado (3,16\%). La ocupación en la actividad registró una tasa anualizada de 2,50\%, por encima del agregado (2,09\%). El PIB por trabajador registró una tasa anualizada de 1,96\%, por encima del agregado (1,05\%). Las horas semanales por trabajador registraron una tasa anualizada de -0,04\%, una caída menos pronunciada que la del agregado (-0,45\%). El PIB por hora trabajada registró una tasa anualizada de 2,00\%, por encima del agregado (1,51\%).

\textbf{En niveles de 2025, el tamaño relativo y la productividad de la actividad económica también importan.} La actividad representó 16,2\% del PIB real agregado, con 165,3 billones de pesos de 2015, y concentró 12,3\% de los ocupados, con 2,8 millones de personas. El PIB por trabajador fue 58,4 millones de pesos de 2015 por ocupado, por encima del agregado (44,5 millones). Las horas semanales por trabajador fueron 43,0, muy cerca del agregado (43,7). El PIB por hora trabajada fue 26,1 mil pesos de 2015 por hora, por encima del agregado (19,6 mil).

\textbf{Zoom a 25 agrupaciones.} Dentro de adm. pública, educación y salud, la apertura a 25 agrupaciones permite ver diferencias que el promedio de la actividad oculta.

\begin{table}[H]
\centering
\caption{Adm. pública, educación y salud: apertura de productividad laboral a 25 agrupaciones, 2010--2025}
\label{tab:sector_o_p_q_zoom25}
\scriptsize
\begin{tabular}{p{0.40\textwidth}rrrr}
\toprule
Actividad económica & PIB/trab. 2025 & Crec. & PIB/hora 2025 & Crec. \\
\midrule
Salud y servicios sociales & 49,1 & 2,65\% & 21,3 & 2,58\% \\
Educación & 52,3 & 1,54\% & 25,8 & 1,26\% \\
Administración pública & 77,2 & 1,40\% & 32,2 & 2,06\% \\
\bottomrule
\end{tabular}
\caption*{\footnotesize Nota: PIB por trabajador en millones de pesos constantes de 2015; PIB por hora en miles de pesos constantes de 2015. Fuente: cálculos propios con DANE y GEIH.}
\end{table}

En esta apertura, el mayor crecimiento del PIB por trabajador se observa en Salud y servicios sociales (2,65\%), mientras que el menor se registra en Administración pública (1,40\%). Por hora trabajada, el mejor desempeño corresponde a Salud y servicios sociales (2,58\%) y el más rezagado a Educación (1,26\%).

\subsection{Artes, otros servicios y hogares}

Esta agrupación reúne actividades artísticas, de entretenimiento y recreación y otras actividades de servicios; actividades de los hogares individuales en calidad de empleadores; actividades no diferenciadas de los hogares individuales como productores de bienes y servicios para uso propio. En la CIIU Rev. 4 A.C. corresponde a R, S y T; divisiones 90--98.

\begin{table}[H]
\centering
\caption{Artes, otros servicios y hogares: PIB, ocupados, horas y productividad laboral, 2010--2025}
\label{tab:sector_r_s_t_productividad}
\small
\begin{tabular}{lrrr}
\toprule
Indicador & 2010 & 2025 & Crec. anual \\
\midrule
PIB real  (Billones de pesos de 2015) & 15,3 & 47,2 & 7,81\% \\
Ocupados (Millones) & 1,5 & 2,0 & 1,91\% \\
PIB por trabajador (Millones de pesos de 2015) & 10,4 & 24,2 & 5,79\% \\
Horas semanales por trabajador & 42,3 & 39,9 & -0,38\% \\
PIB por hora trabajada (Miles de pesos de 2015) & 4,7 & 11,6 & 6,20\% \\
\bottomrule
\end{tabular}
\end{table}

\textbf{Frente al agregado nacional, la comparación variable por variable muestra diferencias relevantes.} El PIB real de la actividad registró una tasa anualizada de 7,81\%, por encima del agregado (3,16\%). La ocupación en la actividad registró una tasa anualizada de 1,91\%, muy cerca del agregado (2,09\%). El PIB por trabajador registró una tasa anualizada de 5,79\%, por encima del agregado (1,05\%). Las horas semanales por trabajador registraron una tasa anualizada de -0,38\%, muy cerca de la caída agregada (-0,45\%). El PIB por hora trabajada registró una tasa anualizada de 6,20\%, por encima del agregado (1,51\%).

\textbf{En niveles de 2025, el tamaño relativo y la productividad de la actividad económica también importan.} La actividad representó 4,6\% del PIB real agregado, con 47,2 billones de pesos de 2015, y concentró 8,5\% de los ocupados, con 2,0 millones de personas. El PIB por trabajador fue 24,2 millones de pesos de 2015 por ocupado, por debajo del agregado (44,5 millones). Las horas semanales por trabajador fueron 39,9, por debajo del agregado (43,7). El PIB por hora trabajada fue 11,6 mil pesos de 2015 por hora, por debajo del agregado (19,6 mil).

\textbf{Zoom a 25 agrupaciones.} Dentro de artes, otros servicios y hogares, la apertura a 25 agrupaciones permite ver diferencias que el promedio de la actividad oculta.

\begin{table}[H]
\centering
\caption{Artes, otros servicios y hogares: apertura de productividad laboral a 25 agrupaciones, 2010--2025}
\label{tab:sector_r_s_t_zoom25}
\scriptsize
\begin{tabular}{p{0.40\textwidth}rrrr}
\toprule
Actividad económica & PIB/trab. 2025 & Crec. & PIB/hora 2025 & Crec. \\
\midrule
Artes, entretenimiento y otros servicios & 32,8 & 5,95\% & 16,0 & 5,89\% \\
Hogares como empleadores & 9,4 & 2,63\% & 4,5 & 3,49\% \\
\bottomrule
\end{tabular}
\caption*{\footnotesize Nota: PIB por trabajador en millones de pesos constantes de 2015; PIB por hora en miles de pesos constantes de 2015. Fuente: cálculos propios con DANE y GEIH.}
\end{table}

En esta apertura, el mayor crecimiento del PIB por trabajador se observa en Artes, entretenimiento y otros servicios (5,95\%), mientras que el menor se registra en Hogares como empleadores (2,63\%). Por hora trabajada, el mejor desempeño corresponde a Artes, entretenimiento y otros servicios (5,89\%) y el más rezagado a Hogares como empleadores (3,49\%).

{\footnotesize Nota general: para facilitar la lectura, el informe usa PIB por actividad económica; en sentido estricto, el numerador corresponde al valor agregado bruto de cada actividad reportado por el DANE. Valores en pesos constantes de 2015. Ocupados expandidos con el factor \texttt{fex}. Las horas semanales promedio se calculan como horas anuales totales divididas por ocupados y por 52 semanas. Se excluye 2020 por no contar con GEIH anual comparable en la base del proyecto. Fuente: cálculos propios con DANE y GEIH.}